\subsection{无穷乘积}

对赋范代数 A 中的点的每个序列 \((\pmb{x}_n)\) ，使它对应于部分积的序列 \(p_n = \prod_{k=0}^{n} x_k\) ；则序列偶 \((x_{n})\) 与 \((p_{n})\) 称为以 \(x_{n}\) 为通项因子的无穷乘积。以 \(x_{n}\) 为通项因子的无穷乘积称为收敛的，如果序列 \((p_{n})\) 在 A 中收敛；此序列的极限则称为序列 \((x_{n})\) 的积，记作 \(\bigoplus_{n=0}^{\infty} x_{n}\) 。

命题 2. 设 \((\pmb{x}_n)\) 是完备赋范代数 A 中的点的一个序列。a) 若以 \(\pmb{x}_n\) 为通项因子的无穷乘积收敛，且 \(\mathbf{P}_{n=0}^{\infty}\pmb{x}_n\) 是 A 中的单位，则对每个 \(\varepsilon > 0\) 存在一个整数 \(n_0\) ，使得

\[
\left\| \prod_ {k = m} ^ {n} x _ {k} - e \right\| \leqslant \varepsilon
\]

每当 \(n_0 \leqslant m \leqslant n\) 。

\begin{enumerate}
\def\labelenumi{\alph{enumi})}
\setcounter{enumi}{1}
\tightlist
\item
  反之，若序列 \((\pmb{x}_n)\) 满足此条件，则以 \(\pmb{x}_n\) 为通项因子的无穷乘积收敛；且若每个 \(\pmb{x}_n\) 都是 A 中的单位，则 \(\mathbf{P}_{n=0}^{\infty}\pmb{x}_n\) 是单位。
\end{enumerate}

此命题的证明逐步仿照定理 1 的证明，留给读者（定理 1 证明中 N 的有限子集 L 应换成区间）。

推论 1. 若以 \(x_{n}\) 为通项因子的无穷乘积收敛，且 \(\bigcap_{n=0}^{\infty}x_{n}\) 是单位，则 \(\lim_{n\to\infty}x_{n}=e\) 。

推论 2. 若以 \(\pmb{x}_n\) 为通项因子的无穷乘积收敛，且 \(\mathop{\sum }\limits_{{n = 0}}^{\infty}{\pmb{x}}_{n}\) 是单位，则以

\[
\mathbf {y} _ {n} = \mathbf {x} _ {n + h} \quad (n \geqslant 0)
\]

为通项因子的无穷乘积收敛。

序列 \((\mathbf{y}_n)\) 的积记作 \(\bigoplus_{n = h}^{\infty}\mathbf{x}_n\) ，也称为无穷乘积的指标 \(h\) 的剩余，该无穷乘积以 \(\mathbf{x}_n\) 为通项因子。

仍在 \(\mathbf{P}_{n=h} x_n\) 是单位的假设下，由命题 2 得：若 \((z_n)\) 是一个序列，使得除去有限个指标外 \(z_n = x_n\) ，则以 \(z_n\) 为通项因子的乘积收敛。

命题 3. 设 \((k_{n})\) 是整数的一个严格递增序列，其项 \(\geqslant 0\) ，且满足 \(k_{0} = 0\) ；若以 \(x_{n}\) 为通项因子的无穷乘积收敛，且

若令

\[
\boldsymbol {u} _ {n} = \prod_ {p = k _ {n}} ^ {k _ {n + 1} - 1} \boldsymbol {x} _ {p},
\]

则以 \(u_{n}\) 为通项因子的无穷乘积收敛，且我们有

\[
\underset {n = 0} {\overset {\infty} {\mathsf {P}}} u _ {n} = \underset {n = 0} {\overset {\infty} {\mathsf {P}}} x _ {n}.
\]

因为序列 \((\boldsymbol{u}_{n})\) 的部分积序列是序列 \((\boldsymbol{x}_{n})\) 的部分积序列的一个子序列。

最后，对阿贝尔群用过的同样论证（第 III 章，§ 5，第 7 号）表明，若一个序列 \((\pmb{x}_n)\) （赋范代数 \(\mathbf{A}\) 中的）可乘，则以 \(\pmb{x}_n\) 为通项因子的乘积收敛，且

\[
\mathop {\text { P }} _ {n = 0} ^ {\infty} x _ {n} = \prod_ {n \in \mathbf {N}} x _ {n}
\]

（它也写作 \(\prod_{n=0}^{\infty} x_n\) ）；反之当然不真（参看习题 7）。
% semantic-fragments: legacy-input-seam
\relax{}
\ExerciseGroupsStart
